\begin{textsample}{POS Dim 3 – 2015\_c6 – Score 1.00 – 2015\_c6/enade\_2015\_ci\_econ\_5389f4fcd9.txt}  \label{ex:f3_pos_215}
CAPÍTULO 3 ANÁLISE TÉCNICA DA PROVA Este capítulo tem por objetivo apresentar o desempenho dos estudantes concluintes de Ciências Econômicas no Enade/2015.
Para isso, foram calculadas as estatísticas básicas da prova em seu todo ( seção 3.1.1 ), bem como as estatísticas relacionadas aos Componentes de Formação Geral ( seção 3.1.2 ) e de Conhecimento Específico da Área ( seção 3.1.3 ).
Dadas as suas características, foram analisadas, em separado, as questões objetivas ( seção 3.2 ) e as questões discursivas ( seção 3.3 ).
Para as questões objetivas, foram disponibilizados os índices de facilidade e de discriminação ponto-bisserial, também em separado, para os Componentes de Formação Geral ( seção 3.2.1 ) e de Conhecimento Específico ( seção 3.2.2 ).
De cada componente, uma das questões foi escolhida para exemplificar a análise gráfica, relacionando as alternativas escolhidas pelos estudantes ( inclusive o gabarito ) com o número de acertos no componente.
O Anexo I apresenta a íntegra da análise gráfica para todas as questões objetivas.
Para cada uma das questões discursivas, os conteúdos dos tipos mais comuns de respostas dos estudantes são apresentados e comparados com o padrão de respostas esperado ( ver Anexo VIII com o padrão de respostas ).
Tomando como base as duas questões discursivas do Componente de Formação Geral, a seção 3.3.1.6 apresenta comentários sobre a correção das respostas com respeito à Língua Portuguesa.
Nas Tabelas 3.1 e 3.3 são apresentados o tamanho da população inscrita, ausente, presente e percentual de faltosos, respectivamente, por Grande Região e por Categoria Administrativa e Organização Acadêmica da IES.
As demais tabelas apresentam as seguintes estatísticas das notas15 : média do desempenho na prova, erro padrão da média, desvio padrão, nota mínima, mediana e nota máxima para cada um de seus componentes.
Tais estatísticas contemplam o total de estudantes concluintes da Área de Ciências Econômicas inscritos e presentes à prova do Enade/2015, tendo em vista agregações ou por Grandes Regiões e o país como um todo, ou por Categoria Administrativa e Organização Acadêmica da IES.
Em relação aos gráficos de distribuição de notas, o intervalo considerado foi de 10 unidades, aberto à esquerda e fechado à direita, com exceção do primeiro intervalo, [ 0 ; 10 ], fechado em ambos os extremos.
Para os gráficos de distribuição das notas das 15 Uma definição dessas estatísticas pode ser encontrada no Glossário. 3.1 CONDIÇÃO DE PRESENÇA E ESTATÍSTICAS BÁSICAS DA PROVA Esta seção apresenta a condição de presença, além das estatísticas selecionadas e histogramas da nota geral ( 3.1.1 ) e de cada componente : Formação Geral ( 3.1.2 ) e Conhecimento Específico ( 3.1.3 ).
São também apresentadas estatísticas selecionadas de subpopulações caracterizadas por Grande Região, Categoria Administrativa e Organização Acadêmica. 3.1.1 Condição de Presença e Estatísticas Básicas Gerais A Tabela 3.1 apresenta a Condição de Presença na prova, por Grande Região, dos estudantes concluintes de Ciências Econômicas.
A população total de Inscritos foi de 9.538.
Destes, 7.853 estiveram Presentes, sendo 17,7 \% o índice de não comparecimento.
A Região de maior abstenção foi a Norte ( 25,0 \%, ) e a de menor abstenção foi a Sudeste ( 16,3 \% ).
Já a Tabela 3.2 apresenta as Estatísticas Básicas da prova, por Grande Região, dos estudantes concluintes de Ciências Econômicas.
A Média das notas da prova como um todo ( nas seções seguintes serão analisados os Componentes de Formação Geral e de Conhecimento Específico ) foi 42,1, sendo que os alunos da região Norte obtiveram a Média mais baixa ( 35,7 ), e os da região Sudeste obtiveram a Média mais alta ( 44,4 ).
As demais médias foram : 37,3 na região Nordeste ; 42,5 na região Sul ; e 42,9 na região Centro-Oeste.
O Desvio padrão para o Brasil como um todo foi 13,7, sendo o maior Desvio padrão encontrado na região Sudeste ( 14,2 ) e o menor, na região Norte ( 10,3 ), indicando uma dispersão um pouco menor das notas desta última região.
A região que obteve a maior nota Máxima foi a Sudeste ( 87,6 ), ao passo que a região que atingiu a menor nota Máxima foi a Norte ( 65,3 ).
A Mediana do Brasil como um todo foi 41,7, sendo a maior Mediana obtida na região Sudeste ( 44,4 ), e a menor obtida no Norte ( 35,4 ).
A nota Mínima foi zero em todas as regiões.
Considerando -se as notas segundo Grande Região, observa -se que existe diferença estatisticamente significativa ao nível de 95 \% entre a menor Média, obtida na região Norte ( 35,7 ), e as médias de todas as demais regiões.
O comportamento das notas dos estudantes de todo o Brasil pode ser observado no Gráfico 3.1 que apresenta um histograma com a distribuição das mesmas.
Essa é uma distribuição unimodal com a moda no intervalo ( 40 ; 50 ].
A Tabela 3.3 apresenta informações referentes à Condição de Presença desagregadas por Categoria Administrativa e por Organização Acadêmica.
Da população total de Inscritos, 3.804 são de IES Privada e 5.734 de IES Públicas.
Em relação à Organização Acadêmica, a maior participação foi obtida por estudantes de Universidades ( 7.888 ), seguido de Faculdades ( 1.033 ) e Centros Universitários ( 617 ).
A Categoria Administrativa de maior abstenção foi a Privada ( 20,9 \% ), e entre as Organizações Acadêmicas foi a dos Centros Universitários ( 27,2 \% ), que ficou acima da abstenção nacional de 17,7 \%.
Já a Tabela 3.4 apresenta informações referentes à Média da nota final desagregadas por Categoria Administrativa e por Organização Acadêmica.
A Média das notas da prova como um todo foi 42,1.
Em relação à Categoria Administrativa, os estudantes das IES Públicas obtiveram Média mais baixa ( 42,0 ), e os das IES Privadas obtiveram Média mais alta ( 42,3 ), que a Média nacional.
Observa -se que não existe diferença estatisticamente significativa entre as médias das notas das IES Públicas e Privadas.
A diferença entre as médias das regiões Sudeste e Norte ( 8,7 ), a maior e a menor Média, foi superior à diferença entre IES Públicas e Privadas ( 0,3 ), caracterizando uma maior diversidade regional do que administrativa.
No tocante à Organização Acadêmica, apenas as Faculdades obtiveram Média mais alta ( 44,4 ) que a nacional.
As demais médias foram : 41,9 nas Universidades e 41,7 nos Centros Universitários.
O Desvio padrão para as IES Públicas ( 13,8 ) e para as Faculdades ( 15,3 ) foi superior ao do Brasil como um todo ( 13,7 ), indicando uma dispersão um pouco maior das notas nesta Categoria Administrativa e nesta Organização Acadêmica.
Constata -se que existe diferença estatisticamente significativa ao nível de 95 \% nas médias das notas dos estudantes de Faculdades em relação aos estudantes de Universidades e Centro Universitários. 3.1.2 Estatísticas Básicas no Componente de Formação Geral A Tabela 3.5 apresenta as Estatísticas Básicas em relação ao componente da prova que avalia a Formação Geral dos estudantes concluintes.
Os alunos de todo Brasil obtiveram desempenho médio de 58,5.
Quanto à variabilidade, o Desvio padrão das notas dos estudantes do Brasil, como um todo, foi 17,0.
A maior Média foi obtida na região Centro-Oeste ( 60,2 ), e a menor, na região Norte ( 55,3 ).
As demais médias foram : 55,9 na região Nordeste ; 59,4 na região Sudeste ; e 59,1 na região Sul.
Já o maior Desvio padrão foi obtido na região Sudeste ( 17,5 ), e o menor, na região Sul ( 16,0 ).
Os demais desvios padrões foram : 16,2 na região Norte, e 16,6 das regiões Nordeste e Centro-Oeste.
A maior nota no Componente de Formação Geral da prova do Enade/2015 foi 98,2, obtida por, pelo menos, um aluno na região Sudeste.
A menor nota Máxima foi obtida na região Norte ( 92,4 ).
A Mediana do Brasil, como um todo, foi 60,2, sendo a menor Mediana encontrada na região Norte ( 55,3 ) ; e a maior encontrada na região Sudeste ( 61,9 ).
A nota Mínima nesta parte foi zero em todas as regiões, sem exceção.
Considerando -se as notas segundo Grande Região, observa -se que existe diferença estatisticamente significativa entre a maior Média das notas do Componente de Formação Geral, obtida na região Centro-Oeste ( 60,2 ), e na região Norte ( 55,3 ), a menor Média.
O Gráfico 3.2 propicia a avaliação do desempenho dos estudantes do Componente de Formação Geral, a partir do histograma da distribuição das notas correspondentes.
A distribuição é unimodal, com moda em ( 60 ; 70 ], enquanto na prova, como um todo ( Gráfico 3.1 ), a moda foi alcançada dois intervalos abaixo ( 40 ; 50 ].
Nota -se, ainda que, no Gráfico 3.2, as notas apresentam maior dispersão do que no Gráfico 3.1 ( distribuição das notas da prova ), confirmado pela comparação dos desvios padrões : 13,7 para a nota da prova como um todo e 17,0 para o Componente de Formação Geral.
Na Tabela 3.6 são apresentadas as informações referentes ao desempenho dos concluintes do Componente de Formação Geral, em diferentes agregações : Categoria Administrativa e Organização Acadêmica.
Observa -se que não existe diferença estatisticamente significativa entre as médias dos tipos de Categoria Administrativa.
A maior Média foi obtida por estudantes de IES Públicas ( 58,9 ), enquanto a menor foi recebida por estudantes de IES Privadas ( 58,0 ).
Considerando -se o tipo de Organização Acadêmica, nota -se que não há diferença estatisticamente significativa entre as médias dos tipos de Organização Acadêmica. 3.1.3 Estatísticas Básicas do Componente de Conhecimento Específico A Tabela 3.7 apresenta as Estatísticas Básicas referentes ao Componente de Conhecimento Específico da área de Ciências Econômicas.
A Média do desempenho dos alunos do Brasil, como um todo, foi 36,6.
A maior Média foi obtida na região Sudeste ( 39,4 ), e a menor, na região Norte ( 29,2 ).
As demais médias foram : 31,1 na região Nordeste ; 37,0 na região Sul ; e 37,1 na região Centro-Oeste.
Quanto à variabilidade das notas, o Desvio padrão do Brasil, como um todo, foi 15,1, sendo o maior Desvio padrão observado na região Sudeste ( 15,6 ), e o menor, na região Norte ( 11,1 ).
Os demais desvios foram : 13,2 na região Nordeste, 14,4 na região Sul e 15,5 na região Centro-Oeste.
A Mediana das notas dos estudantes de todo o Brasil foi 35,6.
A maior Mediana ocorreu na região Sudeste ( 39,0 ), e a menor, na região Norte ( 28,8 ).
As demais medianas foram : 29,6 na região Nordeste ; 37,0 na região Sul ; e 36,1 na região Centro-Oeste.
A nota Máxima do Brasil, como um todo, foi 89,1, sendo obtida por, pelo menos, um aluno na região Sudeste.
As demais notas máximas foram : 63,9 na região Norte ; 77,6 na região Nordeste ; 87,6 na região Sul ; e 86,1 na região Centro-Oeste.
A nota Mínima foi zero em todas as regiões.
Observa -se que existe diferença estatisticamente significativa da menor Média das notas do Componente de Conhecimento Específico da região Norte ( 29,2 ) em relação às demais regiões.
O intervalo de confiança da região Norte ( 2,0 ) é significativamente maior do que de todas as outras.
Assim como os Gráficos 3.1 e 3.2, o Gráfico 3.3, apresentado a seguir, permite uma avaliação do desempenho de concluintes de Ciências Econômicas em relação ao Componente de Conhecimento Específico com um histograma da distribuição das notas correspondentes.
Dentre as três distribuições apresentadas, esta é a mais concentrada nas notas baixas.
Esta também é uma distribuição bimodal, e os grupos modais são o ( 20 ; 30 ] e ( 30,40 ].
A Tabela 3.8 apresenta uma comparação dos resultados em relação à Categoria Administrativa e à Organização Acadêmica, agora levando em conta o desempenho de alunos do Componente de Conhecimento Específico da prova.
No que se refere à Organização Acadêmica, a maior Média foi das Faculdades ( 39,7 ), vindo a seguir a das Universidades ( 36,3 ) e, depois, a dos Centros Universitários ( 36,2 ).
O maior Desvio padrão, e acima do valor para o Brasil, como um todo, foi o das Faculdades ( 17,3 ).
Além de receberam a maior Média, as Faculdades também obtiveram as maiores notas Máxima ( 89,1 ) e Mediana ( 39,0 ).
Por sua vez, as Universidades e os Centros Universitários obtiveram nota Máxima 87,6 e 79,0, respectivamente.
As medianas foram 35,2, para as Universidades, e 35,8, para os Centros Universitários.
A nota Mínima foi zero para as Universidades, 7,4 para os Centro Universitários e 3,7 para as Faculdades.
Quanto à Categoria Administrativa, observa -se um comportamento semelhante àquele da parte de Formação Geral e à prova como um todo, ou seja, não existe diferença estatisticamente significativa entre as médias das IES Públicas ( 36,4 ) e IES Privadas ( 37,1 ).
Neste caso, a maior Média foi obtida por alunos de IES Privadas de ensino.
Observa -se que existem diferenças estatisticamente significativas ao nível de 95 \% no Componente de Conhecimento Específico entre as notas das Faculdades em relação aos Centros Universitários e Universidades. 3.2 ANÁLISE DAS QUESTÕES OBJETIVAS Esta seção apresenta estatísticas selecionadas e histogramas das Questões Objetivas de Formação Geral ( 3.2.1 ).
São também apresentadas e comparadas as médias das subpopulações caracterizadas por Grande Região. 3.2.1 Componente de Formação Geral A Tabela 3.9 apresenta as Estatísticas Básicas relativas às oito questões objetivas do componente da prova que abrange a Formação Geral dos estudantes.
A Média do Brasil foi 61,1.
A menor Média foi encontrada na região Norte ( 55,9 ), e a maior, na região Centro-Oeste ( 63,9 ).
As demais médias foram : 56,9 na região Nordeste ; 62,5 na região Sudeste ; 62,1 na região Sul.
O Desvio padrão do Brasil foi 20,8, sendo o maior Desvio padrão encontrado na região Nordeste ( 21,1 ), e o menor, na região Norte ( 19,8 ).
Os demais desvios foram : 20,9 na região Sudeste ; 19,9 na região Sul ; e 20,6 na região Centro-Oeste.
As medianas ( 62,5 ), as notas máximas ( 100,0 ) e as notas mínimas ( 0,0 ) foram iguais para todas as regiões.
A Tabela 3.10 apresenta o Índice de Facilidade e o Índice de Discriminação ( ponto-bisserial ) para cada uma das questões objetivas do Componente de Formação Geral.
Quanto ao Índice de Facilidade, foram usadas as seguintes cores para diferenciar o nível de dificuldade da questão : • Azul para as questões classificadas com índice Muito fácil ( \&gt ; =0,86 ), verde para as questões classificadas com índice Fácil ( 0,61 a 0,85 ), amarelo para as questões classificadas com Médio ( 0,41 a 0,60 ), vermelho para as questões classificadas com Difícil ( 0,16 a 0,40 ) e roxo para as questões classificadas com Muito difícil ( \&lt ; =0,15 ).
Já quanto ao Índice de Discriminação, foram usadas as seguintes cores para qualificar a questão : • As questões classificadas com índice Fraco receberam a cor vermelho ( \&lt ; =0,19 ), as classificadas com Médio receberam a cor amarelo ( 0,20 a 0,29 ), as classificadas com Bom receberam a cor verde ( 0,30 a 0,39 ) e as classificadas com Muito bom ( \&gt ; =0,40 ) receberam a cor azul.
As questões objetivas do Componente de Formação Geral, segundo o Índice de Facilidade, foram assim avaliadas : das oito questões, nenhuma teve o Índice de Facilidade classificado como Muito fácil.
Duas questões foram tidas como Fácil, por terem índice de acertos situado na faixa entre 0,61 e 0,85 ( de 61,0 \% a 85,0 \% de acertos ).
Quatro questões foram consideradas de dificuldade Médio, situando -se no intervalo entre 0,41 e 0,60 do Índice de Facilidade, ou seja, houve entre 41,0 \% e 60,0 \% de acertos, enquanto duas questões foram classificadas na categoria Difícil, situando -se no intervalo entre 0,16 e 0,40.
Por fim, nenhuma das questões apresentou menos de 15 \% de acertos, razão pela qual seria classificada como Muito difícil.
Como já comentado, para análise das questões objetivas, relativas à Formação Geral, segundo o poder de discriminação, utilizou -se o Índice de Discriminação ( ponto- A Tabela 3.10 apresenta o Índice de Facilidade e o Índice de Discriminação ( ponto-bisserial ) para cada uma das questões objetivas do Componente de Formação Geral.
Quanto ao Índice de Facilidade, foram usadas as seguintes cores para diferenciar o nível de dificuldade da questão :  Azul para as questões classificadas com índice Muito fácil ( \&gt ; =0,86 ), verde para as questões classificadas com índice Fácil ( 0,61 a 0,85 ), amarelo para as questões classificadas com Médio ( 0,41 a 0,60 ), vermelho para as questões classificadas com Difícil ( 0,16 a 0,40 ) e roxo para as questões classificadas com Muito difícil ( \&lt ; =0,15 ).
Já quanto ao Índice de Discriminação, foram usadas as seguintes cores para qualificar a questão :  As questões classificadas com índice Fraco receberam a cor vermelho ( \&lt ; =0,19 ), as classificadas com Médio receberam a cor amarelo ( 0,20 a 0,29 ), as classificadas com Bom receberam a cor verde ( 0,30 a 0,39 ) e as classificadas com Muito bom ( \&gt ; =0,40 ) receberam a cor azul.
As questões objetivas do Componente de Formação Geral, segundo o Índice de Facilidade, foram assim avaliadas : das oito questões, nenhuma teve o Índice de Facilidade classificado como Muito fácil.
Duas questões foram tidas como Fácil, por terem índice de acertos situado na faixa entre 0,61 e 0,85 ( de 61,0 \% a 85,0 \% de acertos ).
Quatro questões foram consideradas de dificuldade Médio, situando -se no intervalo entre 0,41 e 0,60 do Índice de Facilidade, ou seja, houve entre 41,0 \% e 60,0 \% de acertos, enquanto duas questões foram classificadas na categoria Difícil, situando -se no intervalo entre 0,16 e 0,40.
Por fim, nenhuma das questões apresentou menos de 15 \% de acertos, razão pela qual seria classificada como Muito difícil.
Como já comentado, para análise das questões objetivas, relativas à Formação Geral, segundo o poder de discriminação, utilizou -se o Índice de Discriminação ( ponto-bisserial ). bisserial ).
Nesta análise, as questões foram assim avaliadas : seis das oito questões apresentaram índice acima ou igual a 0,40 e, assim, foram classificadas com índice Muito bom para esse grupo de estudantes.
Duas questões tiveram Índice de Discriminação Bom, com valor entre 0,30 e 0,39, para esse grupo de estudantes.
Nenhuma questão teve nível Médio ou Fraco de discriminação para esse grupo de estudantes.
O Índice de Facilidade variou de 0,35 a 0,85, e o de Discriminação, de 0,36 a 0,55.
As seis questões com Índice de Discriminação Muito bom figuraram entre níveis de dificuldade Fácil e Médio, desse conjunto : cinco classificadas na categoria Fácil ( questões 2, 3, 5, 6 e7 ) do Índice de Facilidade e uma, na categoria Médio ( questão 8 ).
Em particular, a questão 7 foi a que apresentou o maior poder discriminatório, com índice 0,55, porém foi considerada fácil em termos de facilidade, com uma proporção de 0,62 acertos.
O máximo de acertos foi alcançado pela questão 5 com um Índice de Facilidade de 0,85.
A questão de número 8 apresentou um Índice de Facilidade de 0,56, ou seja, mais da metade dos estudantes conseguiu resolvê -la, dentro do universo de participantes.
Seu Índice de Discriminação foi Muito bom ( 0,48 ).
O Gráfico 3.4, para exemplificar, analisa o comportamento da questão de número 2 de Formação Geral.
Trata -se de uma questão considerada mediana, com relação à facilidade e a que obteve o maior Índice de Discriminação dessa parte da prova.
Neste gráfico, cada uma das cinco curvas representa o percentual de respostas em determinada alternativa da questão, em função do número de acertos dos estudantes nessa parte da prova ( Formação Geral/Múltipla Escolha ), antes de possíveis eliminações pelo critério do ponto bisserial.
A curva em preto corresponde à alternativa C, a correta para esta questão.
Assim, observa -se que entre os estudantes com menor número de acertos, nessa parte do exame, a situação mais frequente foi a escolha da alternativa D ( em roxo ), incorreta.
Por exemplo, entre os estudantes que acertaram três questões, 33,7 \% escolheram a alternativa D, 28,9 \% a alternativa C ( correta ), 13,0 \% a alternativa A, 14,0 \% a alternativa B e 9,8 \% a alternativa E.
Entre os que acertaram três respostas entre as questões de múltipla escolha de formação geral, 0,5 \% deixou a questão em branco e 0,2 \% marcou mais de uma alternativa, invalidando a questão.
À medida em que o número de acertos aumenta, indicando desempenho melhor nessa parte da prova, aumenta concomitantemente a proporção de estudantes que selecionaram a alternativa correta C, atingindo 100 \% para os estudantes com 8 acertos.
Essa análise permite verificar como a questão discriminou os grupos de desempenho, justificando o alto índice obtido na questão.
Cumpre notar que não é possível inferir deste gráfico nem o índice de facilidade ( que seria uma média da proporção ponderada pela quantidade de alunos com cada uma das notas ), nem o índice de discriminação ponto bisserial, por razão equivalente.
No caso extremo no qual a grande concentração dos acertos dos alunos fosse abaixo de 4, o índice de facilidade seria obrigatoriamente abaixo de 40 \% ( neste exemplo ).
Caso a concentração fosse em 6 acertos ou mais, o índice seria obrigatoriamente acima de 60 \%.
Os gráficos relativos às demais questões de Formação Geral constam do Anexo I. 3.2.2 Componente de Conhecimento Específico A Tabela 3.11 apresenta as Estatísticas Básicas em relação às questões objetivas do Componente de Conhecimento Específico da prova de Ciências Econômicas por Grande Região.
A Média do Brasil deste componente foi de 40,1.
A menor Média foi observada na região Norte ( 32,5 ), e a maior, na região Sudeste ( 42,8 ).
O Desvio padrão de todo o Brasil foi 16,3, sendo o menor Desvio padrão encontrado na região Norte ( 12,2 ), e o maior, na região Sudeste ( 16,9 ).
A Mediana de todo o Brasil foi 39,1, a mesma encontrada nas regiões Sul e Centro-Oeste.
As outras regiões apresentaram os seguintes valores para a Mediana : 30,4 na região Norte ; 34,8 na região Nordeste ; e 43,5 na região Sudeste.
A nota Máxima da prova foi 95,7, obtida nas questões objetivas do Componente de Conhecimento Específico, por, pelo menos, um aluno na região Sudeste.
A menor nota Máxima foi obtida na região Norte foi 69,6, enquanto a maior foi verificada na região Sudeste ( 95,7 \% ) e nas demais regiões, a nota Máxima da prova foi 91,3.
Em todas as regiões a nota Mínima foi zero.
A Tabela 3.12 apresenta os Índices de Facilidade e Discriminação ( ponto-bisserial ) das questões objetivas do Componente de Conhecimento Específico para os estudantes de Ciências Econômicas.
Para facilitar a diferenciação das questões, usaremos as mesmas cores da Tabela 3.10 para as diferentes classificações dos Índices de Facilidade e de Discriminação.
Dentre as questões objetivas da parte da prova relativa ao Componente de Conhecimento Específico, nenhuma questão foi anulada pela Comissão.
Desse modo, a classificação, quanto ao Índice de Facilidade, foi estabelecida com base nas 27 questões.
A partir dos índices obtidos, pode -se concluir que a maioria das questões objetivas da prova foi considerada, pelo menos, Difícil : das 27 questões válidas, 17 foram classificadas como Difícil ou como Muito difícil ( classificação modal ).
Três questões foram classificadas como Fácil ou como Muito fácil, e outras sete consideradas como Médio.
Já quanto aos índices de discriminação das questões objetivas do Componente de Conhecimento Específico da prova, tem -se como resultado a seguinte classificação : nove das 27 questões foram consideradas como boas, enquanto quatro delas tiveram Índice de Discriminação Muito bom.
Assim, para quase metade das questões – 13 em 27 – os Índices de Discriminação foram Bom ou Muito bom.
Dentre as demais, dez delas foram classificadas como Médio e outras quatro como Fraco, sendo 14, por conseguinte, a quantidade de questões nos dois patamares mais baixos de discriminação.
Constata -se, assim, que a prova – no que se refere ao Componente de Conhecimento Específico – possuía boa capacidade de discriminar entre aqueles que dominam ou não o conteúdo.
Dentre as questões que alcançaram os maiores índices de discriminação, nove delas, as de números 12, 13, 14, 19, 21, 24, 26, 28 e 34, foram classificadas com Índice de Discriminação Muito bom, situando -se no intervalo de 0,40 a 0,50 do índice, três delas ( questões 12, 28 e 34 ) foram classificadas na categoria Fácil, quanto ao Índice de Facilidade, sete ( 13, 14, 19, 24, 26, 27 e 30 ) na categoria Médio e 16 ( questões 9, 11, 15, 16, 17, 18, 20, 21, 22, 23, 25, 29, 31, 32, 33 e 35 ), na categoria Difícil.
A questão de número 10 foi a mais difícil dentre as 27 questões específicas válidas, com baixo Índice de Facilidade, apenas 9,0 \% de acertos.
Destacam -se, as questões 9, 11, 18 e 31, com Índice de Discriminação baixo, igual a 0,16 para cada questão.
Tais questões foram, portanto, pelo critério ponto-bisserial, consideradas inadequadas.
Por isso, as questões 9, 11, 18 e 31 foram eliminadas do cômputo da nota final.
A título de exemplo das análises do comportamento das questões objetivas, o Gráfico 3.5 analisa a questão 19 do Componente de Conhecimento Específico.
Apesar de esta questão ter sido considerada pelos estudantes avaliados como uma questão de classificação Médio como nível de facilidade da prova, pois apresentou Índice de Facilidade 0,55, ou seja, 55,0 \% dos estudantes assinalaram acertadamente a opção E, correspondente ao gabarito.
Seu Índice de Discriminação foi igual a 0,49, classificado como Muito bom, também sendo esta questão a que apresentou o maior índice discriminatório.
Neste gráfico, cada uma das cinco curvas representa o percentual de respostas em determinada alternativa da questão 19, em função do número de acertos dos estudantes nessa parte da prova, antes de possíveis eliminações de questões pelo critério do ponto-bisserial.
A alternativa correta E, representada no gráfico pela curva em vermelho, foi escolhida em maiores proporções pelos alunos com desempenho melhor nessa parte da prova.
Já as alternativas incorretas, também denominadas distratores, foram selecionadas, principalmente, por aqueles com notas mais baixas.
Neste caso também a soma não é sempre 100 \% por causa das questões não respondidas ou com mais de uma opção marcada.
Aqueles com nota zero, na sua quase totalidade deixaram esta questão em branco ou marcaram mais de uma alternativa, comportamento considerado inválido.
A proporção de alunos que selecionou a resposta correta E aumenta gradativamente, chegando a atingir 100 \% para 21 acertos ou mais, enquanto a proporção dos que escolheram alternativas incorretas decai, a partir de três acertos, como função do número de acertos nessa parte da prova.
Os gráficos relativos às demais questões do Conhecimento Específico constam do Anexo I. 3.3 ANÁLISE DAS QUESTÕES DISCURSIVAS Esta seção apresenta estatísticas selecionadas e histogramas das Questões Discursivas de Formação Geral ( 3.3.1 ) e Conhecimento Específico ( 3.3.2 ), são também apresentadas e comparadas as médias de subpopulações caracterizadas por Grande Região. 3.3.1 Componente de Formação Geral As análises dos resultados de desempenho dos estudantes de Ciências Econômicas nas duas questões discursivas relativas à Formação Geral encontram -se na Tabela 3.13 e no Gráfico 3.6.
Na Tabela 3.13, observa -se que a nota Média nesse conjunto de questões foi próxima da obtida nas objetivas.
Os estudantes, de todo o Brasil, obtiveram, em Formação Geral, Média 61,1 nas questões objetivas e 54,7 nas questões discursivas.
No entanto, pode -se notar um aumento do Desvio padrão de 20,8 nas questões objetivas do Componente de Formação Geral dos alunos de todo o Brasil, para 25,1 nas questões discursivas do mesmo componente.
A maior Média foi obtida na região Sudeste ( 54,8 ), e a menor, na região Norte ( 54,3 ).
A Mediana de todo o Brasil, neste componente, foi 62,5.
Nas regiões Norte, Nordeste e Sul a Mediana foi 62,0.
A Mediana foi maior na região Sudeste ( 63,5 ) e menor na região Centro-Oeste ( 61,8 ).
A nota Máxima da prova foi 98,0, obtida por, pelo menos, um aluno na região Sudeste.
Nas demais regiões a nota Máxima foi : 94,0 na região Norte ; 97,0 na região Nordeste ; 96,5 na região Sul ; e 95,0 na região Centro-Oeste.
A nota Mínima ( 0,0 ) foram as mesmas em todas as regiões do Brasil.
O Gráfico 3.6 representa a distribuição das notas nas questões discursivas do Componente de Formação Geral.
A moda desta distribuição ocorre no intervalo ( 60 ; 70 ].
Destaca -se, também, o intervalo [ 0 ; 10 ] com distribuição próxima a 12,5 \% do total de notas, sendo que no intervalo [ 0 ; 10 ] incluem -se além da nota zero, a frequência de alunos que deixaram este tipo de questão em branco.
Na sequência, os resultados verificados para cada uma das questões discursivas de Formação Geral serão apresentados, estabelecendo -se relações com os conteúdos abordados em cada uma delas.
Os comentários da Banca de docentes corretores, a respeito do observado na correção das respostas dos estudantes, suas impressões e conclusões serão apresentados junto à análise de cada questão.
Cumpre esclarecer que, tendo em vista que as questões discursivas de Formação Geral são padronizadas, ou seja, constam de todas as provas, os comentários da Banca são os mesmos para todas as carreiras acadêmicas, sendo direcionados a todos os estudantes que participaram do Enade/2015.
A seguir, serão analisados os desempenhos da Área de Ciências Econômicas nas duas questões discursivas de Formação Geral do Enade/2015, comparando -se os resultados obtidos com comentários para cada questão. 3.3.1.1 Análise de Conteúdo da Questão Discursiva 1 do Componente de Formação Geral Os dados de Ciências Econômicas, obtidos a partir das respostas à questão 1, encontram -se na Tabela 3.14 e no Gráfico 3.7.
Nessa questão – de melhor desempenho dentre as duas de Formação Geral – os alunos, de todo o Brasil, tiveram Média 60,4.
A maior Média para a questão 1 foi obtida na região Centro-Oeste ( 60,9 ), e a menor, na região Norte ( 60,1 ).
Quanto à variabilidade das notas, o Desvio padrão de todo o Brasil foi 27,3.
O menor Desvio padrão foi obtido na região Nordeste ( 26,1 ), e o maior Desvio padrão foi obtido na região Sudeste ( 28,1 ).
A Mediana ( 70,0 ) e a nota Mínima ( 0,0 ) foram as mesmas em todas as regiões do Brasil.
A nota Máxima ( 100,0 ) também foi a mesma para todas as regiões do Brasil, exceto a região Norte que obteve nota Máxima 90,0.
O Gráfico 3.7 mostra a distribuição das notas na questão discursiva 1 do Componente de Formação Geral.
Observa -se que a maior frequência corresponde aos alunos que obtiveram nota dentro do intervalo ( 70 ; 80 ], correspondendo à moda da distribuição.
Destaca -se, também, os alunos que deixaram a questão em branco com quase 15 \% do total em cada uma das situações. 3.3.1.2 Comentários sobre a correção de Conteúdo das respostas à Questão Discursiva 1 O enunciado da Discursiva 1 de Formação Geral estava claro e indicava um tema de extrema relevância, qual seja : a busca pela educação universal, especificamente, que apontava um caso onde ocorreu a discriminação de acesso ao ensino formal atingindo uma menina, representando a categoria das mulheres discriminadas em tantos campos da vida e, mais ainda, em sociedades que cultuam o poder soberano dos homens.
A linguagem utilizada foi clara, indicando a existência dos fatos vinculados à luta da menina Malala pelo acesso ao ensino formal das jovens do seu país, e o confronto ideológico que gerou o ataque à sua integridade física.
Diante da repercussão internacional, houve o reconhecimento pelo prêmio Nobel, e a ONU resolveu apoiá -la como símbolo da luta pela igualdade entre os sexos buscando a universalização da educação.
O conteúdo está integrado às Diretrizes, vez que o acesso ao ensino formal está vinculado ao direito constitucional à educação, que, também, é protegido pelo Estatuto da Criança e do Adolescente e nos pactos de Direitos Humanos aderidos pelo Brasil.
Por outro lado, o texto demandava uma reflexão sobre a realidade dos indivíduos que não têm acesso à educação e a situação das mulheres excluídas desse processo por segmentos da sociedade.
Como se tratava de texto contendo uma ampla gama de possibilidades, era natural a ocorrência de interpretações diversas das esperadas.
Releva notar que o enunciado pretendeu, primordialmente, identificar a questão de gênero, apontando para as dificuldades que as mulheres possuem, não somente no Brasil, mas em vários países do mundo.
Ao analisar as provas realizadas, constatamos que no item ‘ a ’ poucos foram os que identificaram os estatutos normativos indicados no padrão de resposta.
Essa ausência pode ser justificada tendo em vista que, no enunciado indicado, não existia elemento que pudesse ser utilizado pelo concluinte para embasar sua resposta nesse caminho.
As respostas foram mais gerais, identificando que a educação constitui um item fundamental para o desenvolvimento do indivíduo e que as mulheres sofrem mais com as limitações impostas por países com regime patriarcal, que discriminam o sexo feminino e, com frequência, proíbem as mulheres de frequentar a escola.
Deve -se ter em vista que a indicação de resposta que vinculava o texto aos estatutos normativos representou um raciocínio próprio dos estudantes de Direito que, naturalmente, encaminharam suas respostas para a fundamentação nessa linha identificando, não somente a Constituição, mas também leis, regulamentos, etc.
Tal linha de argumentação não foi objeto da formação dos estudantes das demais áreas.
A maior parte dos concluintes estabeleceu o tema educação formal associado à busca pelo conhecimento, como acesso ao ensino, indicando a independência intelectual e o desenvolvimento da consciência crítica, limitando as manipulações e permitindo amplo acesso ao mercado de trabalho, traduzindo uma libertação econômica e ascensão social.
No tocante ao item ‘ b ’, restou claro que a maioria esmagadora das respostas confluiu para as indicações do padrão apresentado, sendo relevante o tema de igualdade de gênero, além de candente para as concluintes do sexo feminino que desenvolveram os textos de forma concatenada, fundamentada e, muitas vezes, lamentando que, nos dias atuais, ainda se discuta a necessidade de igualar os gêneros na sociedade.
Nessa linha, foram identificadas respostas, consoante com o padrão apresentado, especialmente nos seguintes temas :  violência física e psicológica contra a mulher, incluindo a Lei Maria da Penha, no caso específico do Brasil ;  tolerância/intolerância a vestimentas, trajes, comportamentos socialmente estereotipados ;  aspectos socioculturais que impõem à mulher uma condição de submissão na sociedade, tais como a impossibilidade de a mulher manifestar seus desejos e posicionamentos em algumas culturas, entre outros ;  igualdade/desigualdade de gênero, por exemplo no mercado de trabalho, em relação à desigualdade salarial ;  ideais de Liberdade, Igualdade e Fraternidade : referência a esses ideais como possibilidade de equilibrar as relações de poder entre homens e mulheres.
Em muitos casos, também, os concluintes fizeram um liame entre o acesso à educação e a melhora das relações entre homens e mulheres ao permitir que a conclusão de maior grau de educação formal permita o acesso a cargos públicos relevantes, empregos bem remunerados e posições de poder, tanto na área pública quanto na área privada.
Como as opções do segundo item foram mais amplas, o padrão de resposta identificou as linhas de pensamento adotadas pelos concluintes.
A maior parte das respostas identificou, claramente, o problema da relação de gênero e suas consequências na sociedade.
Também foi frequente constatarem que existe efetiva melhora na situação das mulheres no Brasil, mas que tal movimento não ocorre em todo o mundo.
Aliás, o exemplo de Malala conduziu a uma reflexão em âmbito mundial, e os concluintes apresentaram suas análises críticas à situação da mulher em países que não admitem sequer a sua inserção no plano educacional e, muito menos, no mercado de trabalho.
A maior parte dos concluintes apresentou concatenação lógica nas suas respostas e uma minoria, pouco expressiva, não demonstrou esse domínio de desenvolvimento lógico dos textos.
Ressalte -se que essa referência não inclui a 3.3.1.3 Análise de Conteúdo da Questão Discursiva 2 do Componente de Formação Geral A Tabela 3.15 mostra que o desempenho médio dos estudantes na questão discursiva 2 ( média 51,7 ) foi inferior ao obtido na questão discursiva 1 ( média 60,4 ).
A região Sudeste foi aquela cuja Média, nessa questão, foi maior ( 52,3 ), e a de menor Média foi a região Norte ( 50,2 ).
Quanto à variabilidade das notas, o Desvio padrão de todo o Brasil foi 28,1, superior ao obtido na questão discursiva 1 ( 27,3 ).
O maior desvio nessa questão foi obtido na região Norte ( 29,1 ), enquanto o menor foi obtido na região Sul ( 26,9 ).
A Mediana ( 61,0 ) e a nota Mínima ( 0,0 ) foram as mesmas em todas as regiões do Brasil.
A nota Máxima ( 100,0 ) foi obtida, por pelo menos, um estudante da região Sudeste.
Nas demais regiões a nota Máxima foi : 99,0 na região Norte ; e 98,0 nas regiões Nordeste, Sul e Centro-Oeste.
O Gráfico 3.8 mostra a distribuição das notas na questão discursiva 2 do Componente de Formação Geral.
Observa -se que a maior frequência corresponde aos alunos que obtiveram nota dentro do intervalo ( 60 ; 70 ], correspondendo à moda da distribuição.
Destaca -se, também, os alunos que ou deixaram a questão em branco com mais de 15 \% do total.
Nota -se, ainda, que as notas ficaram mais distribuídas ao longo dos intervalos em comparação à questão discursiva de número 1. 3.3.1.4 Comentários sobre a correção de Conteúdo das respostas à Questão Discursiva 2 O enunciado da Discursiva 2 estava, igualmente, claro, indicando um tema de extrema relevância, qual seja : o reconhecimento de manifestações artísticas nascidas na periferia, como integrante da cultura reconhecida pela sociedade.
A linguagem utilizada foi clara, indicando a análise dos fatos e suas consequências no cotidiano dos partícipes e o confronto com outras visões culturais que são consideradas mais adequadas pela sociedade civil.
O conteúdo está integrado às Diretrizes, vez que o aspecto cultural e sua inter-relação com o conhecimento haurido no nível superior é essencial para a compreensão de mundo do universitário concluinte, com reflexo na sua atuação profissional.
O texto demanda uma reflexão sobre a realidade dos indivíduos que não têm acesso aos itens básicos da sociedade, como educação, saúde, cultura tradicional e segurança e que, apesar disso, criam alternativas para o seu lazer.
A incorporação do funk como manifestação cultural e artística da sociedade brasileira encontrou, basicamente, duas correntes de resposta : • os que aquiescem positivamente à integração ; • os que rejeitam, acerbamente, tal possibilidade.
No meio, uma corrente que somente admite a integração com modificações na conduta dos integrantes do movimento, com a exclusão de letras exclusivamente com apelos sexuais e incitação à violência.
A interpretação positiva indica a origem da música na periferia, notadamente nas comunidades cariocas e sua projeção para o Brasil e o mundo.
Nessa perspectiva, surge o reconhecimento como cultura popular comparada ao samba, axé e outros ritmos também nascidos no Brasil.
Também ocorre comparação com o movimento hip-hop americano.
Já a análise negativa prioriza a vinculação do ritmo com traficantes, prostitutas, vilipêndio às mulheres, falta de controle quanto ao som, local dos eventos, transações sexuais em público, uso frequente de drogas, violência exacerbada.
Faltam, para essa corrente, elementos básicos da cultura, tratando -se de música sem autonomia como movimento de mudança ou aspecto próprio.
A corrente intermediária reconhece o vínculo originário com as comunidades dominadas por quadrilhas de traficantes, mas admite o seu reconhecimento de acordo com a mudança dos atores que transitam na modalidade, como músicos, intérpretes, para adequar as letras das canções expurgando os aspectos considerados nocivos como o apelo sexual, as indicações positivas para a violência e apologia às drogas.
Um aspecto bastante criticado nas músicas dessa modalidade foi a exploração da mulher, entendida como um objeto de satisfação masculina.
É relevante tal constatação, pois a questão pertinente à Malala reflete a luta pela valorização da mulher.
Os concluintes compreenderam a ideia que foi problematizada e escolheram uma estratégia adequada para solucioná -la.
O caráter dinâmico da cultura foi comumente expressado, não através do enfrentamento abordado no texto, mas mediante a indicação ( de forma rasa e genérica ) das dificuldades e consequente aceitação social de manifestações culturais, antes marginalizadas e que hoje integram o acervo cultural nacional, tais como : samba, grafite, capoeira e forró.
Quanto às reflexões de Laraia, a grande maioria reconheceu o funk como manifestação cultural, porém, o identificou mais na sua expressão musical, no balanço de seu ritmo, tecendo contundentes críticas, e mesmo forte discordância quanto ao conteúdo e mensagem das letras quando remetem à depreciação da imagem feminina ou induzem à apologia ao crime.
De resto, os textos abordaram bem mais o item dois do padrão de Respostas.
Uma minoria enfrentou os temas mudança e dinamismo cultural.
As respostas da corrente positiva indicam que o preconceito sempre acompanhou as manifestações populares como o samba, na sua origem.
Admitem que diante da miscigenação surgem movimentos culturais de várias naturezas, o que inclui o funk.
Também indicam a dificuldade de lazer dos membros das comunidades o que faz surgirem eventos para preencher essa necessidade.
Nesse âmbito, apontam a música como alternativa de ascensão social e integração aos bens de consumo.
As respostas da corrente negativa vinculam a música ao tráfico de drogas cujos elementos comandam as comunidades, propiciando que os bailes sejam mecanismos de venda e consumo dessa mercadoria ilícita.
Na mesma toada, indicam a vinculação com a exploração sexual, principalmente das mulheres maiores e menores de idade.
Indicam aspectos imorais nas letras.
Apresentam os locais dos bailes como inadequados e provocadores da desordem pública.
Recusam que músicas com discriminação às mulheres sejam consideradas como cultura.
A corrente intermediária indica a possibilidade de adoção do funk como cultura desde que ocorram modificações nas letras e no comportamento dos atores ( músicos, letristas e público ), retirando os aspectos negativos já apontados.
Fala -se em “ funk do bem ”.
A maior parte das respostas apresentadas indica que os concluintes têm conhecimento dos conflitos sociais e, no caso da questão do funk, conseguem identificar a origem popular da música.
A partir daí é que ocorre a cisão entre os que defendem o movimento como cultura, aduzindo a ascensão social dos artistas envolvidos dos que vinculam a música com atos ilícitos, rejeitando sua inserção no âmbito cultural.
Em algumas provas, há dificuldade de concatenação lógica de raciocínio, mas a maior parte consegue desenvolver o texto de forma razoável.
Essa questão permite visualizar a divisão social com os pensamentos liberais e conservadores, sendo bastante claros nas respostas dos concluintes. 3.3.1.5 Análise de Língua Portuguesa das Questões Discursivas do Componente de Formação Geral Os dados de Ciências Econômicas, obtidos a partir das respostas às questões discursivas do Componente de Formação Geral, no que tange à Língua Portuguesa, encontram -se na Tabela 3.16 e no Gráfico 3.9.
Nesse aspecto, os alunos, de todo o Brasil, tiveram Média 54,1.
A maior Média com respeito à Língua Portuguesa foi obtida na região Sudeste ( 54,2 ), e a menor, na região Nordeste ( 53,8 ).
Quanto à variabilidade das notas, o Desvio padrão de todo o Brasil foi 25,4.
O menor Desvio padrão foi obtido na região Sul ( 24,5 ) e o maior Desvio padrão foi obtido na região Sudeste ( 26,1 ).
A Mediana das notas de Língua Portuguesa foi 61,5, obtida na região Norte.
A maior Mediana foi obtida na região Sudeste ( 62,0 ), e a menor, na região Centro-Oeste ( 60,5 ).
Nas regiões Nordeste e Sul a Mediana foi 61,0.
A nota Máxima para todo o Brasil foi de 99,0, com, pelo menos, um aluno tirando essa nota na região Sudeste.
Nas demais regiões a nota Máxima foi : 96,5 na região Norte ; 98,0 na região Nordeste ; 98,5 na região Sul ; e 97,5 na região Centro-Oeste.
Além disso, a nota Mínima foi zero em todas as regiões do país.
O Gráfico 3.9 mostra a distribuição das notas de Língua Portuguesa do Componente de Formação Geral.
Observa -se que a maior frequência corresponde aos alunos que obtiveram nota no intervalo ( 60 ; 70 ]. 3.3.1.6 Comentários sobre a correção das respostas de Formação Geral com respeito à Língua Portuguesa Ao encaminhar as questões 1 e 2 na direção da produção de um texto dissertativo, espera -se que o participante utilize seus conhecimentos sobre o assunto e estruture seus textos de acordo com as exigências do registro formal próprio dessa situação comunicativa.
Essa configuração determina exigências quanto : à adequação da seleção vocabular, ao desenvolvimento do conteúdo, à estruturação sintática dos períodos, à organização lógica das ideias, à utilização de procedimentos de encadeamento textual e referenciação, à obediência às exigências morfossintáticas próprias da modalidade escrita da norma-padrão, ao respeito às regras ortográficas e às regras de acentuação gráfica.
O padrão de resposta utilizado na avaliação das questões 1 e 2 considerou os aspectos relevantes ao bom desempenho linguístico como competências distintas, de modo a permitir um mapeamento detalhado do domínio dos recursos disponíveis na Língua Portuguesa para a comunicação escrita formal.
Com base nesse objetivo, foram avaliados os seguintes aspectos : a ) Estruturação textual condizente com o gênero solicitado ( texto dissertativo ) e o modo de organização textual expositivo adequado ao gênero – essa competência envolve : • distribuição do conteúdo do texto em parágrafos, de modo a garantir a sua organização temática ; • estruturação sintática condizente com o padrão da modalidade escrita formal da língua portuguesa de modo a garantir a clareza necessária ; • utilização de operadores discursivos que contribuam para a progressão temática do texto, estabelecendo relações lógicas entre as ideias apresentadas, tanto do ponto de vista intrafrasal, como do interfrasal ; • utilização de procedimentos de referenciação lexical e pronominal que permitam a retomada de referentes textuais ; • utilização de sinais de pontuação que contribuam para a organização lógica da frase e do texto.
Espera -se, portanto, que o participante recorra a procedimentos linguístico-discursivos para organizar seu texto, permitindo o encadeamento lógico entre suas partes de forma a garantir a progressão e a coerência textual.
Isso significa que os seguintes procedimentos foram penalizados, de acordo com o padrão de resposta proposto : • elaboração de frases fragmentadas que comprometam a estrutura lógico-gramatical do texto ; • sequência justaposta de ideias sem encaixamentos sintáticos, reproduzindo hábitos da oralidade ; • elaboração de frase com apenas oração subordinada, sem oração principal ; • emprego equivocado do conector ( preposição, conjunção, pronome relativo, alguns advérbios e locuções adverbiais ) que não expresse a relação lógica adequada entre dois trechos do texto e prejudique a compreensão da mensagem ; • repetição ou substituição inadequada de palavras sem se valer dos recursos oferecidos pela língua ( pronome, advérbio, artigo, sinônimo ) ; • utilização inadequada dos sinais de pontuação comprometendo a clareza textual. b ) Respeito às convenções ortográficas da norma-padrão da Língua Portuguesa – essa competência envolve o domínio das regras de acentuação gráfica e da grafia padrão das palavras ( com ausência de abreviaturas próprias da linguagem da internet ), de acordo com as convenções estabelecidas pela legislação em vigor e consubstanciadas no Vocabulário Ortográfico da Língua Portuguesa, editado pela Academia Brasileira de Letras ( com aceitação da legislação anterior, no caso das regras relativas ao uso do hífen e da acentuação gráfica, já que o exame foi realizado no mês de novembro de 2015 ).
A competência relativa ao respeito às convenções ortográficas envolve :  grafia correta das palavras ;  respeito às regras de acentuação gráfica ;  emprego de maiúsculas em início de frase, em nomes próprios de pessoas, lugares ou instituições ;  ausência de abreviações como p/, vc, tb, pq, tá, né, usadas muitas vezes em escrita informal e na internet ;  obediência às regras de separação de sílabas no final da linha. c ) Domínio dos diferentes aspectos morfossintáticos próprios da modalidade escrita formal da norma-padrão da Língua Portuguesa – os itens relacionados a essa competência são : a concordância nominal, a concordância verbal, a regência nominal, a regência verbal, a flexão nominal, a flexão verbal, a correlação entre os tempos verbais, a colocação pronominal.
O domínio dessa competência envolve :  flexão do verbo para estabelecer concordância de número com o sujeito da frase ;  flexão do artigo, do adjetivo e do pronome para concordar em número e em gênero com o substantivo a que se referem ;  respeito à regência nominal e à verbal, utilizando a preposição adequada depois de um substantivo, de um verbo ou de um adjetivo ; emprego da preposição antecedendo o pronome relativo e emprego adequadamente os pronomes relativos “ cujo(a) ” e “ onde ” ;  emprego adequado do acento grave indicador da crase entre uma preposição e um artigo ( a+a ) ;  atendimento às regras de colocação pronominal ( próclise e ênclise ), distintas dos hábitos da oralidade ou da escrita informal. d ) Seleção vocabular adequada à modalidade escrita formal da Língua Portuguesa, exigida pela situação comunicativa – essa competência envolve a precisão na utilização do vocabulário relacionado à temática solicitada pela questão ; a ausência de marcas da oralidade, como termos de sentido muito genérico ( “ coisa ”, “ negócio ”, “ você ” ) e termos de registros mais informais ( como gírias, jargões, frases feitas, ditados populares, termos regionais ).
Assim, espera -se que o participante respeite a adequação vocabular não empregando gírias ou expressões coloquiais, evite repetição desnecessária de palavras e utilize um vocabulário mais formal, como solicitado por um texto dissertativo. ( “ coisa ”, “ negócio ”, “ você ” ) e termos de registros mais informais ( como gírias, jargões, frases feitas, ditados populares, termos regionais ).
Assim, espera -se que o participante respeite a adequação vocabular não empregando gírias ou expressões coloquiais, evite repetição desnecessária de palavras e utilize um vocabulário mais formal, como solicitado por um texto dissertativo.
A escolha dessas competências para subsidiar o processo de avaliação apoia -se na concepção de que, no desempenho dos graduandos, a modalidade escrita tem apresentado uma intensa simplificação, originada no padrão da modalidade oral da Língua Portuguesa.
No caso do texto de base dissertativa, inscrito em um registro formal, a distância entre as duas modalidades é ainda maior, o que provoca situações de hipercorreção ( desvios provocados pela utilização inadequada de uma regra da norma-padrão ) e de truncamentos sintáticos ( estruturas frasais incompreensíveis devido à complexidade sintática própria da modalidade escrita ).
Observam -se, então, os seguintes aspectos que marcam essa distinção entre as duas modalidades, devido à excessiva simplificação da modalidade falada : a ) redução drástica de estruturas subordinadas, compensada pelo aumento na frequência de estruturas coordenadas e absolutas, por um lado, ou pela elaboração de estruturas truncadas pelo excesso de ideias sem a devida conexão subordinativa ; b ) redução no uso de operadores argumentativas para expressar relações lógicas essenciais à construção do texto, substituídas pela exigência de inferência por parte do interlocutor para suprir a sua ausência ; c ) redução do uso do subjuntivo, ao lado da ampliação do uso do indicativo combinado a estruturas frasais coordenadas ou absolutas ; d ) empobrecimento do processo de referenciação, com a repetição de pronomes ou nomes ; e ) simplificação extrema da marcação da categoria tempo na morfologia verbal ; f ) falta de domínio de vocabulário mais abstrato e de maior complexidade, essencial ao desenvolvimento do processo dissertativo ; g ) redução drástica no emprego da acentuação gráfica, processo intensificado pela intensa utilização das redes sociais.
Os aspectos macroestruturais da elaboração do texto não foram avaliados neste processo porque dizem respeito à avaliação do conteúdo.
Portanto, a banca de Formação Geral, composta por profissionais de diferentes áreas do conhecimento, se encarregou -se da avaliação do atendimento ao solicitado no enunciado das questões do ponto de vista do desenvolvimento do conteúdo.
A avaliação do desempenho linguístico considerou, portanto, três grandes grupos de competências, segundo os aspectos explicitados anteriormente : a ) domínio das convenções ortográficas : grafia de vogais e consoantes, uso de maiúsculas e minúsculas, emprego do hífen e acentuação gráfica ; b ) domínio dos procedimentos de estruturação textual do ponto de vista microestrutural : organização interna dos períodos, emprego de conectores para a articulação lógica entre os períodos e entre os parágrafos, emprego de marcas de referenciação lexical e pronominal ; utilização dos sinais de pontuação que contribuem para a organização lógica da frase. c ) domínio das regras de caráter morfossintático estabelecidas como modelares do ponto de vista da modalidade escrita formal da norma-padrão da Língua Portuguesa : concordância nominal e verbal, regência nominal e verbal, colocação pronominal, flexão nominal e verbal, correlação entre tempos e modos verbais, ausência de marcas de oralidade.
A seleção vocabular adequada à modalidade escrita formal da Língua Portuguesa foi incorporada à essa última competência, tendo em vista a intersecção entre as duas do ponto de vista das exigências do registro formal da modalidade escrita da norma-padrão.
Apresenta -se, a seguir, o padrão de resposta aprovado pelo Inep e já utilizado na avaliação do desempenho linguístico das questões 1 e 2 de Formação Geral no Enade/2013 e no Enade/2014.
Às competências, reunidas nos três grupos descritos anteriormente, foram atribuídos os seguintes pesos relativos : aspectos ortográficos ( 20 \% ) ; aspectos textuais ( 40 \% ), aspectos morfossintáticos e vocabulares ( 40 \% ).
Com base na avaliação do processo desenvolvido nos dois anos anteriores, utilizaram -se quatro níveis para a avaliação do desempenho linguístico nas três competências.
O nível zero ficou reservado para casos especiais, como fuga total ao tema, palavras soltas, produção de uma frase incompleta, entre outros.
A avaliação das respostas elaboradas pelos participantes revelou um resultado coerente na comparação entre as duas questões.
Entretanto, em relação aos resultados do Enade/2014, observou -se uma melhora do domínio de todas as competências, talvez devido ao perfil dos cursos analisados, como Direito e Jornalismo, que pressupõem a utilização profissional da linguagem.
Observou -se uma diferença no desempenho dos participantes em relação às duas questões discursivas, provavelmente em função dos temas solicitados : - a Questão 1 solicitava que o participante desenvolvesse um texto dissertativo sobre o significado da premiação de Malala, abordando dois aspectos : o direito das jovens à educação formal e as relações de poder entre homens e mulheres no mundo.
Esse encaminhamento favoreceu o desenvolvimento mais articulado do texto, o que se refletiu em melhor desempenho dos aspectos vocabulares e textuais. - a Questão 2 solicitava que o participante desenvolvesse um texto dissertativo para discutir a questão do reconhecimento do funk como legítima manifestação artística e cultural do povo brasileiro.
Por ser um tema muito polêmico, o tema do preconceito sócio-cultural-econômico favoreceu a construção de textos mais críticos, dividindo o conjunto de participantes em dois grupos : os que defendiam esse tipo de manifestação artística como legítima e os que a criticavam de maneira veemente, chegando, inclusive, a utilizar muitos termos de baixo calão.
Essa última atitude teve como consequência a elaboração de textos mal desenvolvidos do ponto de vista sintático-discursivo, com muitos truncamentos.
Os resultados revelaram essa diferença.
Ao analisar os dados de cada questão, constata -se que a Questão 1 teve maior número de questões de médio ou bom desempenho, em relação à Questão 2.
Por outro lado, a Questão 2 apresentou um índice maior de notas baixas em relação à Questão 1.
Os resultados da avaliação correspondem aos seguintes aspectos observados em cada competência : a ) aspectos ortográficos : o desempenho dos participantes revelou uma diferença muito grande nos dois aspectos analisados : baixo índice de desvios ortográficos e grande índice de desvios de acentuação.
Em vários casos, ocorre ausência completa de acentuação gráfica.
Os resultados revelam que a tendência dominante entre os universitários brasileiros é a eliminação da acentuação gráfica, talvez motivada pelos hábitos relacionados às redes sociais e pela ausência de esclarecimento dos meios de comunicação, das autoridades e das escolas sobre as decisões do Acordo Ortográfico da Língua Portuguesa de 1990.
Os casos mais sistemáticos de eliminação do acento indicador da sílaba tônica são :  palavras proparoxítonas ( “ proximos ”, “ politicos ” ) ;  palavras paroxítonas terminadas em ditongo crescente ( “ necessario ”, “ noticiarios ”, “ individuo ”, “ dependencia ”, “ varios ”, “ propria ” ) ; n palavras oxítonas ( “ ninguem ”, “ esta ”, “ ate ”, “ tambem ”, “ alvara ” ) ; n uso indevido do acento gráfico em palavras como “ genêro ”, “ melâncolia ”, por exemplo.nnQuanto ao domínio das convenções relativas à grafia das palavras, observam -se desvios como : nI. hipercorreção pela escolha de “ e ” no lugar de “ i ” ou de “ i ” no lugar de “ e ”, por influência de hábitos da oralidade ( “ descriminação ” por “ discriminação ”, “ entervenção ” por “ intervenção ”, “ esteriotipada ” por “ estereotipada ” ) ; nII. eliminação do “ r ” marcador do infinitivo verbal ( “ esta ” no lugar de “ estar ”, “ estuda ” no lugar de “ estudar ” ). nIII. eliminação de uma sílaba na palavra “ educação ”, que foi grafada “ edução ” por um grande número de participantes. nIV. outros casos de desvios de grafia relacionados à variação diastrática podem ser observados em “ estrupo ”, “ subjulgadas ”, “ precoseito ”, “ soubre ”, “ vecendo ”, “ indesencia ”, “ apolojia ”, “ fulga ”. nV. desvios de segmentação podem ser atribuídos à confusão entre o morfema gramatical “ mos ” e o pronome oblíquo “ nos ”, em “ chegar mos ”, “ enfrentar-mos”.nVI. inadequação no uso da maiúscula para destacar determinadas palavras-chave do texto, como “ Violência ”, “ Brasileiros ”.
Destaque -se, também, o grande número de participantes que grafam os textos inteiramente em caixa alta. nnVale observar, também, que, ao contrário do que se esperava, não apareceram abreviaturas próprias do “ internetês ”, relacionado ao uso de redes sociais e emails.nnb ) aspectos textuais : esta competência é a que se revela como a mais problemática entre os participantes, porque são muitos os problemas observados, desvios acumulados durante toda a formação do estudante : sequência justaposta de ideias sem encaixamentos sintáticos ; redução drástica de estruturas subordinadas, ao lado do aumento na frequência de estruturas coordenadas e absolutas ; redução no uso de conectores para expressar relações lógicas essenciais à construção do texto, substituídas pela exigência de inferência por parte do interlocutor para suprir a sua ausência ; emprego equivocado de operadores que não estabelecem relações lógicas coerentes entre ideias do texto ; emprego inadequado do pronome relativo ( com omissão da preposição ou a utilização de pronome inadequado, como “ onde ” ) ; repetição exaustiva de termos sem a utilização de procedimentos mais sofisticados de substituição ( hiperonímias, hiponímias, nominalizações, expressões metafóricas ) ; frases fragmentadas que comprometem a estrutura lógico-gramatical ; frases formadas apenas por oração subordinada, sem oração principal. a utilização de pronome inadequado, como “ onde ” ) ; repetição exaustiva de termos sem a utilização de procedimentos mais sofisticados de substituição ( hiperonímias, hiponímias, nominalizações, expressões metafóricas ) ; frases fragmentadas que comprometem a estrutura lógico-gramatical ; frases formadas apenas por oração subordinada, sem oração principal.
Um importante aspecto a destacar é o baixíssimo desempenho de uma parte dos participantes em relação à estrutura formal do texto produzido, o que é extremamente preocupante ao se levar em conta que são graduandos em fase final de formação.
São frequentes os casos de desvios de estruturação frasal, com uso inadequado ou ausência de conectivos entre parágrafos e entre frases.
Em uma parte dos textos, falta um mínimo de textualidade e de domínio do registro padrão da língua.
Na verdade, observam -se relações linguísticas quase agramaticais, como as estabelecidas pela sequência de gerúndios sem o apoio de um ponto de partida para a organização das informações gramaticais e semânticas.
Observou -se que uma grande parte dos participantes não organizou as ideias em parágrafos, talvez devido ao pequeno número de linhas disponibilizadas para a resposta da questão ou, quem sabe, pela suposição de que não seria necessária essa divisão por não se tratar de um texto no modelo de uma redação dissertativo-argumentativa, como solicitado nos vestibulares.
Quanto à utilização dos mecanismos de referenciação, deve -se destacar a ocorrência, em uma boa parte dos textos, de repetições de palavras ou expressões sem a utilização de termos sinônimos ou pronomes, como seria adequado.
Outro aspecto relevante a ser destacado é a quase total ausência de operadores argumentativos, tanto intrafrasais como interfrasais, repercutindo uma tendência atual da mídia escrita.
Quanto à utilização dos sinais de pontuação, observou -se uma grande precariedade nos textos analisados.
É muito frequente a ocorrência de parágrafos sem marca interna de pontuação para separar os períodos.
Vale observar que não foi penalizada a ausência de vírgula para destacar locuções ou adjuntos adverbiais de pequena extensão deslocados de posição na frase, por ser um uso opcional.
São os seguintes os tipos de problemas encontrados : I. vírgula : utilização de vírgula para separar o sujeito e o predicado ; ocorrência de apenas uma das vírgulas para separar uma palavra, uma expressão ou uma oração encaixada ; uso de vírgula no lugar do ponto para separar ideias que constituem períodos distintos ; II. ponto e vírgula : utilização do ponto e vírgula no lugar de vírgula ; III. ponto final : ausência de ponto final para separar períodos. c ) aspectos morfossintáticos e vocabulares : os resultados são muito transparentes em relação aos aspectos mais problemáticos no desempenho dos participantes.
O desvio mais frequente, em relação à regência, é a falta do sinal indicativo da crase – isso revela que o usuário não tem consciência de que, sob a forma do termo “ a ”, existe a presença de uma preposição “ a ”, exigida pela regência do termo anterior, e de um artigo definido.
Embora em outros exames, como o Enem, a falta de crase seja penalizada em acentuação, nesta avaliação esse desvio foi considerado no âmbito dos aspectos morfossintáticos.
Outro problema relacionado à regência verbal e à nominal, encontrado frequentemente nas questões, foi a ausência de preposição antes de pronome relativo, processo generalizado na modalidade oral da língua, em situações de registro informal.
Apesar da possibilidade de que essa alteração de regência se generalize no padrão escrito da Língua Portuguesa, como já está ocorrendo até em textos jornalísticos, o não emprego da preposição foi penalizado neste processo de avaliação.
Outro desvio muito frequente diz respeito aos processos de concordância verbal e de concordância nominal.
Quanto à concordância de número, observou -se ausência de marca com sujeito posposto ou uso indevido ( uso inadequado da marca de plural comandado pelo núcleo plural da locução adjetiva, apesar de o substantivo que funciona como núcleo do sintagma nominal estar no singular ).
Uma ocorrência generalizada foi a ausência de acento circunflexo na forma plural do presente do indicativo do verbo “ ter ”, que foi considerada como um desvio na concordância verbal e não na acentuação gráfica.
Quanto à concordância de gênero, vários casos foram observados, normalmente no âmbito de sintagmas nominais longos, em que o adjetivo está afastado do substantivo.
Quanto à questão da colocação pronominal, foram poucos os desvios observados.
Apesar de serem aspectos relacionados à oralidade, concluiu -se que, no registro escrito formal, a maioria dos participantes já incorporou regras como a não introdução da frase por um pronome oblíquo e a próclise na presença de um termo atrator.
Não se adotou, entretanto, o padrão excessivamente formal descrito pelas gramáticas normativas em relação à posição do pronome oblíquo em locuções verbais, já que esse uso está muito distante da prática cotidiana, até em textos mais formais.
Quanto aos aspectos vocabulares, alguns tipos de inadequação foram observados, principalmente na Questão 2, relacionada ao funk : expressões da oralidade, seleção vocabular incompatível com o contexto, gerando situações de falta de inteligibilidade ; falta de domínio de vocabulário mais abstrato e de maior complexidade, essencial ao desenvolvimento do texto de base dissertativa.
Ainda em função da temática desenvolvida na Questão 2, foi significativa a presença de palavras de baixa calão, por retratar um universo cultural específico.
Os textos que continham esse tipo de vocabulário foram considerados inadequados em função do registro formal exigido e foram avaliados com grau zero.
Várias marcas de oralidade foram identificadas, embora não em alta frequência : o uso do pronome relativo “ onde ” como relativo universal, falta de artigo definido antes de substantivo, repetição de palavras por falta de vocabulário, reduções como “ tá ” e “ pra ”, expressões informais. 3.3.2 Componente de Conhecimento Específico Na parte da prova relativa às questões discursivas do Componente de Conhecimento Específico ( Tabela 3.17 ), observa -se que a Média foi bem mais baixa do que para as questões discursivas do Componente de Formação Geral.
Enquanto no Componente de Formação Geral a Média para estudantes de Ciências Econômicas de todo o Brasil foi 54,7, na parte de Conhecimento Específico a Média foi 17,3.
A maior Média deste componente foi obtida pelos estudantes da região Sudeste ( 20,1 ), e a menor, pelos da região Norte ( 10,8 ).
Quanto à variabilidade das notas, o Desvio padrão de todo o Brasil foi 16,6.
O maior Desvio padrão foi encontrado na região Sudeste ( 17,5 ), e o menor, na região Nordeste ( 13,0 ).
A maior nota Máxima foi obtida na região Sudeste ( 91,7 ), enquanto a menor nota Máxima foi encontrada na região Norte ( 58,3 ).
Além disso, a nota Mínima ( 0,0 ) foi obtida por pelo menos um aluno em todas as regiões do Brasil.
A Mediana do Brasil como um todo foi 13,3, a mesma obtida na região Centro-Oeste.
A maior Mediana foi obtida na região Sudeste ( 18,3 ), e a menor, na região Norte ( 5,0 ).
Na região Nordeste a Mediana foi 8,3, e na região Sul 14,2.
O Gráfico 3.10 representa a distribuição das notas nas questões discursivas no Componente de Conhecimento Específico.
A moda desta distribuição ocorre no primeiro intervalo, [ 0 ; 10 ].
Nota -se que a distribuição é decrescente.
A análise de cada uma destas questões será feita a seguir. 3.3.2.1 Análise da Questão Discursiva 3 do Componente de Conhecimento Específico Na questão 3, cujos resultados aferidos encontram -se descritos na Tabela 3.18, a Média dos estudantes de todo o Brasil foi 27,2.
O desempenho dos estudantes nessa questão foi superior ao das questões 4 e 5.
A menor Média nessa questão foi obtida pelos alunos da região Norte ( 17,2 ), enquanto a maior Média foi obtida na região Sudeste ( 31,7 ).
Quanto à variabilidade das notas, o Desvio padrão de todo o Brasil foi 30,2.
O maior Desvio padrão foi obtido na região Sudeste ( 31,8 ), enquanto o menor foi obtido na região Norte ( 25,3 ).
A nota Máxima ( 100,0 ) e nota Mínima ( 0,0 ) foi alcançada por, pelo menos, um aluno de todas as regiões do Brasil.
A Mediana do Brasil como um todo foi 20,0, a mesma nas regiões Sudeste e Sul.
Na região Centro-Oeste a Mediana foi 15,0, e na região Nordeste foi 5,0.
Já na região Norte a Mediana foi zero, o que significa que mais de 50 \% dos alunos obteve nota zero ou deixou a questão em branco nessa região do país.
O Gráfico 3.11 mostra a distribuição das notas na questão discursiva 3, do Componente de Conhecimento Específico da área de Ciências Econômicas.
Sem levar em consideração que quase 25 \% dos estudantes deixaram a questão em branco, essa distribuição é unimodal com moda na nota zero. 3.3.2.2 Comentários sobre as respostas à Questão Discursiva 3 A questão foi adequada e compatível com os objetivos do Enade/2015, que avalia alunos ao final da graduação de Ciências Econômicas.
A demanda pela leitura do gráfico, envolvendo habilidades de identificação da tendência e de síntese de indicadores macroeconômicos são fundamentais para a atuação do economista no mercado de trabalho.
Visando a essa formação geral e ao mundo do trabalho, tais conhecimentos são exigidos pelo Ministério da Educação para a conclusão do ensino de Economia.
Aspectos presentes nas Diretrizes do MEC que são avaliados pela Questão 3 : i ) Avaliação de indicadores ; ii ) Análise de tendência ; iii ) Identificação de fenômenos e decisões econômicas e seus impactos.
A questão trouxe um problema com aplicabilidade no mercado de trabalho.
As análises pedidas são abordadas nos períodos iniciais dos cursos de graduação, previstas em qualquer currículo básico do ensino de Ciências Econômicas.
O conhecimento dos agregados macroeconômicos e a avaliação do seu papel no crescimento econômico são, tipicamente, apresentados aos alunos logo nos primeiros períodos da graduação.
A aplicação a dados reais e a necessidade de analisar a conjuntura econômica, com base em extrapolações, tornam a questão mediana em termos de complexidade.
A questão ofereceu os requisitos e informações necessários para sua resolução.
Também o tempo foi considerado suficiente para a redação da resposta.
O padrão de respostas acabou se revelando muito mais complexo e com um espectro de respostas muito mais amplo que o observado na correção.
Basicamente, os alunos que responderam corretamente concentraram suas respostas em dois itens do padrão de resposta.
Os alunos que erraram não o fizeram em divergência ao conteúdo exposto no padrão de respostas.
Em grande medida, erraram a questão aqueles que não interpretaram corretamente o gráfico.
Como destacado no item anterior, o padrão de resposta caracterizou -se por apresentar um amplo espectro de respostas possíveis que poderiam ser aceitas na correção.
Sendo assim, todas as respostas consideradas corretas estavam previstas no padrão de resposta.
No entanto, os alunos que realizaram o exame redigiram respostas que ficaram restritas a apenas alguns dos casos possíveis.
A correção também revelou a compreensão do aluno em relação ao que foi pedido no enunciado, fazendo com que as respostas corretas, apesar de numericamente pequenas, refletissem a proposta da questão e as erradas refletissem falta de conhecimento do tema ou dos instrumentos analíticos necessários à solução do problema.
Como a questão solicitava a interpretação da tabela em consonância com o gráfico, não houve repetição do enunciado nas respostas, fato muito percebido, por exemplo, na Questão 4.
Retomando as Diretrizes do MEC e a aplicação dos conhecimentos adquiridos em cursos de Ciências Econômicas, o elevado número de alunos que não soube interpretar a tendência do gráfico pode ser reflexo de que a linguagem e a abordagem avaliadas não são correspondentes ao que é ensinado na graduação.
O desempenho dos alunos ficou aquém do desejável.
Como já comentado, o grau de complexidade da questão pode ser considerado médio, já que se trata da aplicação de conhecimentos de um rol básico de atribuições do profissional de Economia, que são ensinados nos períodos iniciais da graduação.
Outro ponto que merece destaque foi a grande quantidade de alunos que tratou da questão como pertinente ao ano de aplicação da prova ( 2015 ), e não aos anos que compunham os indicadores selecionados ( 2012 a 2014 ), evidenciando falta de atenção às informações constantes do enunciado.
Notas fracas, em grande medida, mostraram que os alunos interpretaram a tabela de forma incorreta, assim como a tendência mostrada no gráfico e, assim, as análises subsequentes ( itens ‘ b ’ e ‘ c ’ ) acabaram correspondendo a um diagnóstico errado.
Notas medianas, tipicamente, foram de estudantes que realizaram a leitura solicitada no item ‘ a ’ de forma correta, mas responderam de forma incompleta aos itens ‘ b ’ e ‘ c ’.
Quase sempre, apenas citavam as causas e consequências, sem a explicação para os fenômenos.
Notas boas foram respostas que cumpriram, completa ou quase completamente, o que a questão solicitava, ou seja, a leitura correta dos indicadores econômicos, a identificação e a explicação das causas e das consequências do processo de desindustrialização.
Os alunos que responderam à questão focaram suas explicações, para as causas e as consequências do processo, basicamente em cinco pontos : ( i ) variação cambial ; ( ii ) baixa produtividade ; ( iii ) elevado custo da produção ; ( iv ) impacto na balança comercial ; e ( v ) desemprego nos setores industriais.
Além do já comentado erro de leitura dos indicadores, outro erro comum foi a referência ao déficit na balança comercial como uma das consequências do processo, sem especificar as possíveis causas para tal, se redução das exportações ou aumento das importações, e sua relação com a variação cambial.
O elevado número de notas iguais a zero reflete uma deficiência na formação dos alunos, no que se refere à capacidade analítica a partir da identificação e interpretação de indicadores macroeconômicos em análises sobre o desempenho da economia brasileira.
Outras deficiências encontradas foram : inabilidade em interpretação de gráficos, não compreensão de extrapolações e de projeções dos indicadores.
Compondo um dos itens fundamentais das Diretrizes do curso de graduação em Ciências Econômicas, a leitura e avaliação de indicadores econômicos são atributos principais do economista, bem como atividades usualmente ensinadas nos cursos.
São atividades presentes em quase todas as disciplinas do curso, sejam nas matérias de 3.3.2.3 Análise da Questão Discursiva 4 do Componente de Conhecimento Específico A Tabela 3.19 contém as informações relativas à questão 4 do conjunto de questões do Componente de Conhecimento Específico.
A Média geral do Brasil foi 8,5, sendo a menor Média registrada na região Norte ( 5,5 ), e a maior, na região Sudeste ( 10,0 ).
O desempenho dos estudantes de todo o Brasil nesta questão foi mais baixo dentre as três questões discursivas desse componente A nota Máxima ( 100,0 ) foi atingida por, pelo menos, um concluinte nas regiões Sudeste e Sul.
Nas regiões Norte e Nordeste a nota Máxima foi 70,0, e na região Centro-Oeste foi 90,0.
A Mediana em todo o Brasil foi 0,0, o mesmo valor foi obtido em todas as regiões do país.
As notas mínimas foram zero em todas as regiões, o que evidencia um alto percentual de alunos ( pelo menos maior do que 50 \% ) que obtiveram nota zero ou deixaram a questão em branco.
O Gráfico 3.12 representa a distribuição de notas da questão discursiva 4, do Componente de Conhecimento Específico.
Essa distribuição é unimodal com moda na nota zero, seguida de perto pelo intervalo ( 0 ; 10 ].
A maior frequência ocorre na categoria em branco.
Observa -se que uma maior proporção de estudantes deixou essa questão em branco, comparando -se com o ocorrido nas demais questões discursivas, como também o percentual de notas zero é superior nesta questão. 3.3.2.4 Comentários sobre as respostas à Questão Discursiva 4 Dentre as questões discursivas, a Questão 4 foi a que apresentou o maior grau de complexidade e, por isso, pode ser considerada como a mais difícil.
Abordou a Teoria Econômica, neste caso, Microeconômica, como um instrumento para a avaliação de um problema aplicável típico de Finanças, o que a tornou bastante contemporânea e adequada para as demandas do mercado de trabalho do economista.
Além disso, a questão enquadra -se entre as Diretrizes do Ministério da Educação relativas ao Ensino de Economia.
O aluno que respondeu corretamente ao que foi solicitado revelou capacidade de desenvolver raciocínios lógicos e consistentes, de ler e compreender textos econômicos e de utilizar adequadamente a Teoria Econômica.
Tal complexidade foi confirmada pelo desempenho dos alunos, que obtiveram nessa questão a menor média dentre as obtidas nas questões discursivas.
Trazer a necessidade de associação entre Microeconomia e Finanças reflete o caráter educativo do exame, já que evidencia aspectos que os cursos de graduação deveriam valorizar por serem exigências da atuação de um profissional da Área.
Neste sentido, apesar de as questões 3 e 5 também cumprirem esse papel, a Questão 4 vai além dos conhecimentos básicos construídos na graduação, iniciais do curso de Economia, exigindo aplicação dos mesmos em uma demanda contemporânea de trabalho.
Essa associação justifica a avaliação da questão como difícil, dada a realidade dos cursos.
Como já apontado, essa questão está alinhada com as Diretrizes do currículo de Ciências Econômicas, atendendo aos seguintes requisitos : I - desenvolver raciocínios logicamente consistentes ; II - ler e compreender textos econômicos ; III - elaborar pareceres, relatórios, trabalhos e textos na área econômica ; V - utilizar o instrumental econômico para analisar situações históricas concretas ; e VI - utilizar formulações matemáticas e estatísticas na análise dos fenômenos socioeconômicos.
Apesar da complexidade de relacionar dois ramos do ensino de Economia – Microeconomia e Finanças – e de interpretar um problema típico de Governança Corporativa, os comandos dos itens ‘ a ’ e ‘ b ’ eram bastante claros e havia suficiência de informações e disponibilidade de dados para resolver o problema apresentado.
Contudo, o comando do item ‘ c ’ tinha uma elaboração confusa que causou dificuldade para sua interpretação, o que resultou em um elevado número de respostas com nota zero nesse item.
Nos itens ‘ a ’ e ‘ b ’, os alunos revelaram, em grande medida, deficiência na compreensão da teoria, bem como não reconhecimento do que são as falhas de mercado e quais são as formas para mitigar esse tipo de problema econômico.
Por fim, o tempo pareceu adequado para a resolução do problema apresentado na questão.
Dentre as questões discursivas, esta foi a que apresentou o maior número de respostas com grande dissonância do padrão de respostas estabelecido pela Comissão Assessora da Área.
A maioria das respostas incorretas não recebeu nenhuma pontuação por não apresentar minimamente algum dos aspectos presentes no padrão de respostas.
Muitos revelaram não saber o que são “ falhas de mercado ”, conteúdo ensinado nas disciplinas de Microeconomia, e também demonstraram dificuldade de aplicar os conhecimentos ensinados a problemas contemporâneos do mercado de trabalho.
Os alunos que responderam erradamente, muitas vezes, replicavam o que estava escrito no texto que compunha o enunciado.
Com exceções das poucas provas que receberam boas notas, a linguagem adotada pelos alunos refletia, em grande medida, uma profunda falta de conhecimento da teoria, em especial, do que era pedido especificamente – identificar a “ falha de mercado ” concernente ao problema proposto.
A “ falha de mercado ”, uma questão técnica e pertinente à Teoria Econômica, foi, muitas vezes, confundida com problemas de ordem moral e ética.
Essa interpretação revelou outro problema que foi o não reconhecimento de que o objetivo fim de uma firma é a maximização do lucro.
Este conceito é apresentado nas cadeiras iniciais de Microeconomia, tornando -o um conhecimento acessível e \textbf{compreensível} aos alunos logo no início do curso.
A avaliação deste objetivo como moralmente inaceitável revela carências na formação do estudante de Economia.
A própria identificação da maximização do lucro numa abordagem moral reflete que a linguagem utilizada para análises de problemas reais é inadequada, além de dissonante de uma abordagem econômica.
Em grande medida, os alunos utilizaram expressões e/ou trechos dos textos no enunciado e, ainda assim, não utilizavam corretamente os conhecimentos econômicos.
Estas constatações são refletidas no desempenho dos alunos.
O resultado ruim reflete a grande quantidade de provas com notas zero.
Nesse grupo, foi observado que os alunos não compreenderam o objetivo da questão e, muitas vezes, apenas replicaram trechos do enunciado.
Destaca -se também que, se por um lado, o número de notas zero foi muito elevado, por outro, os poucos alunos com notas boas revelaram entender exatamente o que a questão solicitava nos itens ‘ a ’ e ‘ b ’.
Já o item ‘ c ’ foi o de maior incidência de erros.
As abordagens utilizadas podem ser divididas em dois grupos.
Para os alunos que acertaram a questão, grande parte respondeu que a “ falha de mercado ” era a assimetria de informação.
Risco Moral e Agente-Principal apareceram um número menor de vezes.
No item ‘ b ’, também referente ao grupo que acertou a questão, poucos identificaram no programa de remuneração a forma eficaz de mitigar o problema de agência.
Menos ainda foram as respostas em que a regulação, fiscalização e/ou controle figuravam como instrumento para solucionar o conflito de agência.
Muitos alunos utilizaram o termo"controle", podendo indicar que estariam apenas replicando o que constava no enunciado.
Como já enfatizado, o item ‘ c ’ foi aquele com o menor percentual de acerto.
Mesmo os alunos que identificavam a “ falha de mercado ” e o instrumento para mitigar o conflito em análise, raramente alcançaram a pontuação máxima por não acertarem o item ‘ c ’.
Esse fato pode estar associado a duas causas : uma delas pode ser a deficiência na interpretação do comando do item, já que muitos alunos confundiram pontos críticos com análise crítica e acabaram por formular uma hipótese para a baixa aderência das empresas do mercado tradicional a boas práticas de governança corporativa ; outra, que consideramos bem provável, pode estar associada à própria formulação inadequada do comando.
Com relação aos erros, observou -se que muitos alunos não sabem aplicar a Teoria Microeconômica ao problema e não conhecem “ falhas de mercado ”.
Vários copiaram parte dos textos que compunham o enunciado como explicação.
Um exemplo deste tipo de erro foi a utilização da maximização do lucro como falha de mercado.
Muitos alunos demonstraram não ter qualquer conhecimento sobre o que foi perguntado.
Como identificado anteriormente, a alta complexidade da questão pode, em grande medida, ser atribuída à demanda de aplicar conhecimentos de dois ramos distintos do conhecimento – Micro e Finanças – em um problema real do ambiente da Governança Corporativa.
Esta complexidade acabou por revelar a ausência dessa habilidade em um número elevado de respostas. 3.3.2.5 Análise da Questão Discursiva 5 do Componente de Conhecimento Específico A Tabela 3.20 contém as informações relativas à questão 5 do conjunto do Componente de Conhecimento Específico.
A nota Média dos estudantes de todo o Brasil foi 16,1.
A maior Média foi registrada na região Sudeste ( 18,6 ), enquanto a menor Média foi registrada na região Norte ( 9,7 ).
Quanto à variabilidade das notas, o Desvio padrão dos alunos do Brasil, como um todo, foi 20,2.
Enquanto o maior desvio foi encontrado na região Sudeste ( 21,1 ), o menor foi encontrado na região Norte ( 16,2 ).
A nota Mínima ( 0,0 ) e a nota Máxima ( 100,0 ) são as mesmas em todas as regiões.
A Mediana em todo o Brasil foi 5,0, o mesmo valor foi obtido nas regiões Sul e Centro-Oeste.
Na região Sudeste a Mediana foi 10,0 e nas regiões Norte e Nordeste 0,0.
O Gráfico 3.13 representa a distribuição das notas da questão discursiva 5 do Componente de Conhecimento Específico.
Essa distribuição é unimodal com moda no intervalo ( 10 ; 20 ].
Destaca -se, o grande número de estudantes que deixaram esta questão em branco, correspondendo a quase 40 \%. 3.3.2.6 Comentários sobre as respostas à Questão Discursiva 5 Dentre as três questões discursivas, essa foi a de enunciado mais direto.
Em cada item, o comando solicitava de forma direta o ponto sobre o qual o aluno deveria discorrer.
O conteúdo exigido compõe a formação básica do estudante de Economia e, por essa razão, a Questão 5 foi a que apresentou a maior média, apesar de ter sido a questão com o maior número de provas em branco.
Esse fato pode indicar que o tempo tenha sido escasso para solucionar todas as questões discursivas.
O nível de dificuldade foi baixo, apesar de a questão requerer que o aluno diferenciasse correntes teóricas, tal qual sugerido nas Diretrizes do MEC para o ensino de Economia, já que as teorias econômicas abordadas são amplamente ensinadas nos cursos, nos períodos iniciais.
Das questões discursivas, esta foi a que apresentou a relação mais direta entre o padrão de resposta e o que o aluno respondeu.
Dentre as repostas analisadas, os que não alcançaram a pontuação máxima, ou dentro do intervalo de notas boas, revelaram falta de conhecimento da teoria econômica, em especial nos itens ‘ a ’ e ‘ c ’.
No item ‘ a ’, apesar dos alunos demonstrarem algum conhecimento sobre a Lei de Say – solicitada no comando –, não mostraram compreensão sobre os mecanismos de ajuste e retomada do crescimento desta corrente do pensamento econômico.
Já no item ‘ c ’, que tratava das expectativas pela abordagem de Keynes, muitos trataram das expectativas racionais ou adaptativas.
De um modo geral, os alunos que responderam à questão corretamente apresentaram argumentos contidos no padrão de respostas, especialmente, nos itens ‘ a ’ e ‘ b ’.
Os erros típicos encontrados nesses itens foram por incompletude do argumento.
A exceção fica para o item ‘ c ’, no qual se observou que, em grande medida, os concluintes não compreendem as diferentes abordagens sobre expectativas na teoria econômica.
A despeito do maior número de respostas em branco, foi nessa questão que os alunos mais se aproximaram do que era solicitado e do padrão de respostas.
Esse fato pode ter uma dupla interpretação/explicação : a objetividade da questão, sem necessidade de aplicar os conhecimentos a um problema, pois era preciso apenas enunciar conceitos presente na totalidade dos cursos de economia, logo nos períodos iniciais, o que não ocorre no caso da Questão 4.
Nas respostas consideradas fracas, grande parte dos alunos apenas enunciou a Lei de Say e citou a relevância do papel do governo na abordagem keynesiana.
Nas respostas medianas, os alunos enunciaram a Lei de Say, mas não concluíram a análise, destacando os preços e salários flexíveis como mecanismo de ajuste.
Também não explicaram o papel das expectativas na abordagem de Keynes e os mecanismos de retomada do crescimento neste escopo analítico.
Nas respostas consideradas boas, os que não alcançaram a nota máxima foram, em grande medida, aqueles que não responderam ao item ‘ c ’.
No item ‘ a ’, frequentemente, os estudantes não concluíram o raciocínio.
Os alunos reconheciam e enunciavam a lei, mas nem iniciavam a explicação sobre o mecanismo de ajuste da economia, via preços e salários flexíveis.
Mas, como regra geral, as respostas não destoavam do padrão de respostas sugerido.
Comparativamente às outras duas questões discursivas, a Questão 5 foi a que possibilitou que os alunos evidenciassem maior familiaridade com conceitos da teoria econômica.
A despeito da média baixa, os que responderam à questão revelaram entender o que estava sendo perguntado e responderam de modo consistente.
Como já destacado, o que foi identificado foram respostas incompletas e, em função do elevado número de respostas em branco, uma provável falta de tempo para responder às três questões.
A potencial deficiência encontrada foi identificada no item ‘ c ’, no qual os alunos confundiram as perspectivas diferentes na teoria econômica sobre o papel das expectativas. 3.3.3 Considerações Finais De acordo com a definição do Inep, o Enade tem por objetivo aferir “ o desempenho dos estudantes em relação aos conteúdos programáticos previstos nas Diretrizes Curriculares do respectivo curso de graduação, e as habilidades e competências adquiridas em sua formação"( Art. 33-D, Brasil, MEC/Inep.
Portaria Normativa nº 40, de 12/12/2007 ).
Com base nesta definição, esse relatório contém uma série de observações relevantes.
Em complementação, nessa seção, são apresentadas considerações gerais relativas a cada uma das questões discursivas do exame aplicado aos alunos de Ciências Econômicas no Enade/2015.
Questão 3 : como já destacado, o conteúdo da questão foi pertinente à profissão do economista e se encaixou no objetivo do exame evidenciado anteriormente.
Demandou habilidades analíticas e atualização com relação à Economia, neste caso, brasileira.
Apesar da complexidade mediana, a questão revelou uma carência na Questão 3 : como já destacado, o conteúdo da questão foi pertinente à profissão do economista e se encaixou no objetivo do exame evidenciado anteriormente.
Demandou habilidades analíticas e atualização com relação à Economia, neste caso, brasileira.
Apesar da complexidade mediana, a questão revelou uma carência na formação dos alunos em transpor a teoria aprendida ( contas nacionais ) para a identificação dos dados e a análise de tendências.
Também revelou deficiências em interpretação de tabelas e gráficos.
Apesar destas constatações, esta foi a questão com a maior média.
Questão 4 : a questão, em que são associados conhecimentos de Finanças, Microeconomia, e, especificamente, Governança Corporativa, foi adequada, tanto em relação às Diretrizes do MEC quanto aos objetivos do exame, na medida em que avaliou conhecimentos atuais e a capacidade do aluno de aplicá -los a problemas reais.
Apesar do caráter educativo e prospectivo da questão, a média baixa, e, em especial, o baixíssimo nível de acerto do item ‘ c ’, refletem a dissonância entre o objetivo da questão e o preparo dos alunos.
Com relação ao item ‘ c ’, muitos realizaram uma avaliação crítica do conflito de agência, associado aos números da tabela, no lugar de apresentarem os pontos críticos.
Portanto, neste item especificamente, a forma de apresentação da questão pode ter influenciado neste tipo de erro.
Os erros referentes aos itens ‘ a ’ e ‘ b ’, contudo, refletiram, em grande medida, falta de conhecimento e de capacidade de aplicar a teoria em problemas reais ( demanda diferente da Questão 5, em que era pedido para discorrer sobre algumas teorias econômicas de modo direto ).
Os alunos tiveram dificuldade e revelaram não saber o que eram as “ falhas de mercado ” à luz da teoria Microeconômica.
Essa identificação indica a necessidade de repensar o ensino da teoria, que não pode estrar distanciado da aplicação dos conhecimentos.
Questão 5 : a dicotomia entre a baixa complexidade da questão e o elevado número de respostas em branco, além de uma média relativamente baixa, expôs ou a falta de tempo para realização da prova ou falta de conteúdo ou, ainda, a falta de interesse dos alunos em realizar o exame.
Nessa questão, ainda que, em muitas provas, fossem identificados graves erros referentes à teoria, vários alunos responderam corretamente, não atingindo a totalidade dos pontos pela incompletude da resposta.
Em resumo, em que pese o baixo desempenho nas três questões discursivas, a prova refletiu como os conteúdos vêm sendo ensinados nos cursos de graduação de Ciências Econômicas.
Apesar de sabermos que o baixo desempenho pode indicar falta de interesse dos alunos em realizar o exame de modo consistente, ficou evidente a baixa capacidade analítica dos formandos que buscaram responder às questões.
Seja por um ou outro motivo, o desempenho do Enade/2015 é uma excelente ferramenta analítica para avaliar os conteúdos trabalhados em sala, bem como a materialização destes conteúdos em problemas reais e aplicáveis, como será exigido pelo mercado de trabalho do economista.

% matched lemmas: compreensível
\end{textsample}
